\subsection{Garbage-collection pauses in the evaluation loop}
\label{sec:gc}

The power traces of the optimized schedules contain short dips
(\cref{fig:traces}). In the timed window of the full pipeline under the
default governor, each run contains six to seven dips of more than
2\,W below the median power, with onsets at about 4.1, 6.0, 8.3, 11.3,
14.8 and 19.1\,s. The intervals between dips grow from 1.9 to 4.3\,s.
This pattern matches CPython's generational garbage collector. A full
collection runs only when the number of long-lived objects has grown by
25\% since the last one, and every full collection traverses all of
them~\cite{python2024gc}. The evaluation harness, like most COCO
evaluation scripts, keeps every detection as a Python dictionary until
the end of the run. For val2017 these are 732,651 detections, each a
dictionary holding a box list, or about 1.5 million objects tracked by
the collector. Each full collection therefore stalls the main thread
for longer, the GPU drains, and the power drops.

To confirm the cause, a variant was added that performs one collection
before the timed window and then disables the cyclic collector until
the window ends. Reference counting still frees all acyclic objects, so
memory use is unaffected for this workload. The predictions are
byte-identical to the other uint8 rows (\cref{tab:ladder}). No run of
this variant contains a dip. The standard deviation of the sampled
input power falls from 0.88--0.90\,W to 0.16--0.20\,W under the default
governor and from 0.79--0.80\,W to 0.07--0.09\,W with locked clocks.
Throughput rises from 213.8 to 230.2 images/s, so the collector costs
7.1\% of the measured throughput. Energy per image falls from 74.5 to
71.6\,mJ. With locked clocks the cost is 6.0\% (217.9 versus
231.8 images/s).

This gain is not counted as a pipeline improvement. A deployed detector
would pass detections on rather than accumulate them, and its collector
would have little to scan. The effect matters for benchmarking. An
evaluation loop that accumulates results penalizes a fast pipeline more
than a slow one, because the same pauses are a larger share of a
shorter run. The pauses also arrive at irregular, growing intervals,
which inflates tail latency in long runs. All comparisons against the
Ultralytics predictor therefore use the rows with the collector
enabled, which is the more conservative choice for the proposed
pipeline. The rest of the paper follows the same rule.
